\documentclass[10pt,a4paper]{amsart}
\usepackage[margin=19mm]{geometry}
\usepackage{amsmath,amssymb,mathtools}
\usepackage[scr=rsfs,cal=euler]{mathalfa}
\usepackage{xfrac}
\usepackage[unicode,hidelinks,bookmarks=false]{hyperref}
\newcommand{\ie}{i.e.\ }
\newcommand{\cf}{cf.\ }
\newcommand{\resp}{respectively\ }
\newcommand{\Subsection}{\subsection}
\providecommand{\nobreakdash}{\nobreak\hbox{-}}
\setlength{\emergencystretch}{2em}
\multlinegap0pt
\usepackage{bm}
\usepackage{bbm}
\usepackage{dsfont}
\usepackage{xspace}
\usepackage{cancel}
\usepackage{mathtools}
\usepackage{amsthm}
\makeatletter
\@ifundefined{proposition}{\newtheorem{proposition}{Proposition}}{}
\makeatother
\newcommand{\T}{^{\mathrm{T}}}
\newcommand{\cVminus}{\mathcal{V}^-}
\newcommand{\cVplus}{\mathcal{V}^+}
\newcommand{\vA}{\mathbf A}
\newcommand{\vC}{\mathbf C}
\newcommand{\vF}{\mathbf F}
\newcommand{\vI}{\mathbf I}
\newcommand{\vK}{\mathbf K}
\newcommand{\vM}{\mathbf M}
\newcommand{\vO}{\mathbf O}
\newcommand{\vn}{\mathbf n}
\newcommand{\vq}{\mathbf q}
\newcommand{\vw}{\mathbf w}
\newcommand{\vx}{\mathbf x}
\title[Model order reduction of piecewise linear mechanical systems]{Model order reduction of piecewise linear mechanical systems using invariant cones}
\author{A. Yassine Karoui, Remco I. Leine}
\date{Source: arXiv:2601.10241v1 (CC BY 4.0)}
\begin{document}
\maketitle

\noindent\emph{Source and licence.} Model order reduction of piecewise linear mechanical systems using invariant cones. Version arXiv:2601.10241v1 (CC BY 4.0), distributed under the Creative Commons Attribution 4.0 International licence, as recorded in the arXiv licence metadata for this exact version and shown on the abstract page https://arxiv.org/abs/2601.10241v1. Full rights notice: licenses/arxiv-2601.10241-CC-BY-4.0.txt.\par\smallskip

\noindent\textit{Editorial scope.} Counted unit: the Proposition whose source label is \texttt{prop:no\_sliding} in the article (source0.tex lines 587--589, source label \texttt{prop:no\_sliding}), delivered together with the article's own complete original proof of it (source0.tex lines 590--605, one \texttt{proof} environment of 16 lines). No mathematics is written, repaired or re-derived: statement and proof are the article's own text line for line. Required context retained uncounted: eq:eom\_firstorder (lines 549--551); eq:sys\_matrices (lines 553--555). Every definition, display and label of those lines is retained unchanged. Omitted from this package and not redistributed: 1--548, 552--552, 556--586, 606--1342. Nothing inside a retained excerpt is omitted or altered: every retained body is delivered exactly as the article's lines are written, so no omission receipt of any kind accompanies this package. Citations: the retained excerpts carry no citation command at all, so no bibliography entry of the article is transcribed and no reference is redistributed. Reference adaptation: none was needed and none was made; every reference inside the delivered unit resolves inside it, so no unresolved reference or citation is produced. Printed slips kept literal: no printed word, symbol, hypothesis or number of the counted statement or of its proof was altered. Layout and class: the article's own class is \texttt{article} with packages including geometry, babel, fontenc, lmodern, textcomp, latexsym, amsthm, amssymb, amsmath, amsfonts, dsfont, mathrsfs, bbm, bm; the wrapper is the lane's pinned \texttt{amsart} editorial wrapper at 10pt on A4, which restates the theorem-like environments the retained text needs and adds the article's own macro definitions that the retained text needs (\texttt{\textbackslash T}, \texttt{\textbackslash cVminus}, \texttt{\textbackslash cVplus}, \texttt{\textbackslash vA}, \texttt{\textbackslash vC}, \texttt{\textbackslash vF}, \texttt{\textbackslash vI}, \texttt{\textbackslash vK}, \texttt{\textbackslash vM}, \texttt{\textbackslash vO}, \texttt{\textbackslash vn}, \texttt{\textbackslash vq}, \texttt{\textbackslash vw}, \texttt{\textbackslash vx}), with extra wrapper packages (bm, bbm, dsfont, xspace, cancel, mathtools). The delivered numbering is the unit's own. Assets: no retained span contains a figure environment, a graphics-inclusion command, a PDF-page inclusion, a table or a TikZ picture; no image file is copied into this package, the package declares no asset dependency and no image file is redistributed with it. Licence: version arXiv:2601.10241v1 of this article is distributed under the Creative Commons Attribution 4.0 International licence, as recorded in the arXiv licence metadata for this exact version and shown on the abstract page https://arxiv.org/abs/2601.10241v1. A full rights notice is delivered at licenses/arxiv-2601.10241-CC-BY-4.0.txt. Characters: every retained excerpt is pure ASCII, so the delivered engine, XeLaTeX with its default Latin Modern text font, reports no missing character.
\begin{align}\label{eq:eom_firstorder}
	\dot{\vx} = \vF(\vx) = \begin{cases} \vA^- \vx & \text{for } \vx \in \cVminus \\ \vA^+ \vx & \text{for } \vx \in \cVplus \end{cases},
\end{align}

\begin{align}\label{eq:sys_matrices}
	\vA^- = \begin{bmatrix} \vO & \vI \\ -\vM^{-1}\vK & -\vM^{-1}\vC \end{bmatrix}, \quad \vA^+ = \begin{bmatrix} \vO & \vI \\ -\vM^{-1}(\vK + k_n\vw\vw\T) & -\vM^{-1}(\vC + c_n\vw\vw\T) \end{bmatrix}.
\end{align}

\begin{proposition}[Absence of Sliding Modes]\label{prop:no_sliding}
	The homogeneous PWL system defined by Eq. \eqref{eq:eom_firstorder} does not exhibit sliding modes on either of its switching boundaries, $\Sigma_{\alpha}$ or $\Sigma_{\beta}$. Trajectories can only cross these boundaries transversally or, in specific instances where the velocity component normal to the boundary is zero, graze them tangentially.
\end{proposition}

\begin{proof}
	We analyze the behavior at each boundary separately.
	\begin{enumerate}
		\item \textit{Behavior at the position-based switching boundary $\Sigma_{\alpha}$:} The switching function is $h_{\alpha}(\vx) = \vw\T\vq = [ \vw\T, \bm{0}\T ]\vx$, with normal $\vn_{\alpha} = [ \vw\T, \bm{0}\T]\T$. The projection of an arbitrary vector field $\vA^{\pm}\vx$ onto this normal direction is
		\begin{align}
			\vn_{\alpha}\T (\vA^{\pm}\vx) = [ \vw\T, \vO\T] \begin{pmatrix} \dot{\vq} \\ \ddot{\vq}^{\pm} \end{pmatrix} =  \vw\T\dot{\vq}.
		\end{align}
		Crucially, the result is independent of whether the contact or no-contact dynamics are considered. Since the projections are identical, $\vn_{\alpha}\T (\vA^+\vx) = \vn_{\alpha}\T (\vA^-\vx)$, their product can never be negative. The condition for a sliding mode is therefore never met. Trajectories must cross if $ \vw\T \dot{\vq} \neq 0$ or graze $\Sigma_{\alpha}$ if $ \vw\T \dot{\vq} = 0$.
		
		\item \textit{Behavior at the force-based switching boundary $\Sigma_{\beta}$:} The switching function is $h_{\beta}(\vx) = k_n(\vw\T\vq) + c_n(\vw\T\dot{\vq}) = [ k_n\vw\T,  c_n\vw\T]\vx$. For any state $\vx \in \Sigma_{\beta}$, the condition $h_{\beta}(\vx) = 0$ holds by definition. Examining the difference between the two vector fields from Eq.~\eqref{eq:sys_matrices}, we find
		\begin{align}
			(\vA^+ - \vA^-)\vx = \begin{bmatrix} \vO & \vO \\ -\vM^{-1}(k_n\vw\vw\T) & -\vM^{-1}(c_n\vw\vw\T) \end{bmatrix} \vx = \begin{bmatrix} \vO \\ - \vM^{-1}\vw h_{\beta}(\vx) \end{bmatrix}.
		\end{align}
		For any $\vx \in \Sigma_{\beta}$, this difference is identically zero. Therefore, the vector field $\vF(\vx)$ is \textit{continuous} across the boundary $\Sigma_{\beta}$. This continuity trivially excludes the possibility of sliding modes.
	\end{enumerate}
\end{proof}

\end{document}
